\documentclass[11pt,a4paper]{article}
\newcommand{\fichetype}{Exercices}
% Thème A4 commun, librement modifiable.
\usepackage[T1]{fontenc}
\usepackage[utf8]{inputenc}
\usepackage[french]{babel}
\usepackage{lmodern,amsmath,amssymb,mathtools,tikz,pgfplots}
\usepackage[margin=2.2cm,headheight=15pt]{geometry}
\usepackage{xcolor,fancyhdr,enumitem,hyperref,needspace}
\usetikzlibrary{arrows.meta,positioning,calc}
\pgfplotsset{compat=1.18,width=11cm,height=6cm}
\definecolor{Ink}{HTML}{202124}
\definecolor{Blue}{HTML}{356B9E}
\definecolor{Muted}{HTML}{737579}
\color{Ink}
\renewcommand{\familydefault}{\sfdefault}
\setlength{\parindent}{0pt}
\setlength{\parskip}{6pt}
\setlist{itemsep=3pt,topsep=4pt}
\pagestyle{fancy}\fancyhf{}
\fancyhead[L]{\small\color{Muted}MATHS LYCÉE}
\fancyhead[R]{\small\color{Muted}2026--2027}
\fancyfoot[L]{\small\color{Muted}\fichetype}
\fancyfoot[R]{\small\thepage}
\renewcommand{\headrulewidth}{0.3pt}
\hypersetup{colorlinks=true,urlcolor=Blue,linkcolor=Blue}
\newcommand{\R}{\mathbb{R}}
\newcommand{\N}{\mathbb{N}}
\newcommand{\E}{\mathbb{E}}
\newcommand{\V}{\mathrm{V}}
\newcommand{\exercice}[3]{\Needspace{6\baselineskip}\vspace{0.5cm}{\color{Blue}\small\textbf{EXERCICE #1}\hfill #2}\par{\large\bfseries #3}\par}

\begin{document}
{\small\color{Blue}PREMIÈRE SPÉCIALITÉ}\par{\LARGE\bfseries Variations, extrema et optimisation}\par{\large Exercices}\par\vspace{0.2cm}\hrule\vspace{0.4cm}
Les exercices progressent de l’application directe vers la recherche. Détailler les calculs et justifier les réponses.
\par\noindent\begin{minipage}{\linewidth}\exercice{1}{Application}{Une parabole}
Étudier $f(x)=x^2-6x+5$ sur $[0,5]$ et donner ses extrema.
\end{minipage}\par
\par\noindent\begin{minipage}{\linewidth}\exercice{2}{Application}{Un cube}
Étudier les variations de $x^3+2x$ sur $\R$.
\end{minipage}\par
\par\noindent\begin{minipage}{\linewidth}\exercice{3}{Entraînement}{Deux intervalles}
Étudier $1/x$ sur son domaine. Peut-on dire qu’elle est décroissante sur tout son domaine ?
\end{minipage}\par
\par\noindent\begin{minipage}{\linewidth}\exercice{4}{Entraînement}{Rectangle}
Un rectangle a un périmètre de $20$ m. Quelle aire maximale peut-il avoir ?
\end{minipage}\par
\par\noindent\begin{minipage}{\linewidth}\exercice{5}{Application}{Parité}
Étudier la parité de $x^4-3x^2$, $x^3+x$ et $x^2+x+1$.
\end{minipage}\par
\par\noindent\begin{minipage}{\linewidth}\exercice{6}{Entraînement}{Une inégalité}
Montrer que $x+1/x\ge2$ pour $x>0$ par une étude de variations.
\end{minipage}\par
\par\noindent\begin{minipage}{\linewidth}\exercice{7}{Approfondissement}{Production entière}
Un bénéfice fictif est $B(n)=-2n^2+14n-5$ pour un entier $0\le n\le8$. Chercher le maximum.
\end{minipage}\par
\par\noindent\begin{minipage}{\linewidth}\exercice{8}{Approfondissement}{Contre-exemple}
Expliquer pourquoi $f\prime(a)=0$ n’implique pas un extremum. Donner un exemple où un maximum existe au bord avec une dérivée non nulle.
\end{minipage}\par
\end{document}
